\chapter*{Abstract}
\addcontentsline{toc}{chapter}{Abstract}

In this thesis, we present the design of a Retrieval-Augmented Generation (RAG) pipeline integrated into a Visual Studio Code (VSCode) extension, serving as an intelligent assistant that run locally for web developers working with the Next.js framework. The pipeline leverages Milvus \cite{wang2021milvus} as the vector database to efficiently store and retrieve information from Next.js documentation, prepared during the data processing phase. For response generation, the system utilizes Alibaba's Qwen \cite{bai2023qwen}, a Large Language Model (LLM). This combination of technologies ensures the system providing accurate, context-aware assistance. By retrieving relevant information and generating responses grounded in the retrieved data, the proposed design enhances productivity and simplifies the development process within the VSCode environment.

% Retrieval-Augmented Generation (RAG) has emerged as a transformative framework in the field of natural language processing, enabling large language models (LLMs) to leverage external knowledge effectively. This research explores the application of RAG techniques to enhance coding workflows within a specialized Visual Studio Code (VSCode) extension tailored for web development. By integrating state-of-the-art embedding models such as BGE-M3 and robust language generation models like Qwen, our system provides precise and context-aware assistance for developers. Key contributions include a modular architecture that combines vector databases, efficient indexing strategies, and RAG-based augmentation, along with a seamless integration of local LLMs through Ollama, ensuring data privacy and high performance. The research also addresses critical challenges, including managing rapidly evolving frameworks like Next.js and improving interoperability across tools such as Firebase and Tailwind CSS. Evaluation metrics focus on retrieval precision, response relevance, and computational efficiency, demonstrating the system’s potential to revolutionize development workflows. This study highlights the power of RAG in bridging the gap between generative AI and real-world coding applications, paving the way for more intelligent, adaptive, and developer-centric tools.